\chapter{INTRODUCTION}
\label{chap:introduction}

\section{Background of the study}

The rapid expansion of marketplace-based e-commerce has significantly
changed the nature of retail competition. On platforms such as Amazon,
sellers and brand operators compete not only through product quality and
price, but also through organic ranking, advertising visibility, stock
availability, listing presentation, and customer review sentiment. These
signals change continuously and are distributed across multiple parts of
the marketplace interface, making manual monitoring increasingly
inefficient for small sellers and individual analysts.

At the same time, digital marketplaces generate a large volume of
structured and semi-structured data. Category rankings reveal shifts in
competitive visibility; product snapshots contain price, stock, review,
image, and listing-quality signals; and customer reviews provide
unstructured evidence of user satisfaction and dissatisfaction. In
principle, these data sources can support more informed competitive
decisions. In practice, however, their value depends on whether they can
be collected reliably, stored consistently, transformed into defensible
analytical indicators, and delivered to decision-makers in a timely and
interpretable format.

Therefore, the central challenge addressed in this thesis is not merely
data availability, but the conversion of dispersed marketplace
observations into actionable intelligence. This requires an integrated
system that combines scheduled data ingestion, a relational analytical
warehouse, machine-learning and rule-based signal extraction, dashboard
visualisation, and push-based intelligence delivery.

\section{Problem statement}

Despite the growing importance of data-driven competitive intelligence,
many monitoring workflows remain heavily dependent on static Business
Intelligence (BI) dashboards. Such dashboards are useful for exploration,
but they usually follow a pull-based operating model: analysts must open
the interface, select filters, inspect charts, and manually infer which
market events deserve attention. This workflow creates cognitive
friction, particularly when the analyst must monitor many products,
metrics, and categories at the same time.

For Amazon sellers, the operational risk lies in the delay between an
event and its interpretation. A competitor may enter the Top 50 ranking,
increase sponsored visibility, reduce price, lose stock availability, or
modify a product image. Each signal may appear small in isolation, yet
their combination can indicate a meaningful competitive move. Without an
automated mechanism for collecting these signals and synthesising them
into concise alerts, important changes may be detected too late or missed
entirely.

This thesis addresses that practical gap by developing a lightweight
Amazon market-intelligence framework. The implemented system uses Apify
actors for scheduled Amazon data acquisition, stores raw and derived data
in Supabase PostgreSQL tables, runs Python analytical modules through
GitHub Actions, and exposes the resulting intelligence through both
dashboards and Telegram-based OpenClaw agents. Rather than claiming
real-time streaming surveillance, the project focuses on scheduled daily
and weekly monitoring that is technically feasible, cost-aware, and
appropriate for a thesis-scale pilot deployment.

\section{Objectives of the study}

The primary objective of this study is to design, implement, and evaluate
an automated Amazon Marketplace Intelligence Framework for scheduled
competitive monitoring. The framework aims to transform raw marketplace
observations into structured analytical signals, dashboard views, and
concise Telegram briefs.

To achieve this objective, the study pursues the following specific
sub-objectives:
\begin{enumerate}
    \item \textbf{Data acquisition and warehousing:} To build an
    automated Apify-to-Supabase pipeline for collecting Amazon category
    rankings, product-level watchlist snapshots, product images, stock
    status, price information, listing metadata, and customer reviews.
    In the codebase, this objective is primarily implemented through
    \texttt{scripts/ingest\_category.py},
    \texttt{scripts/ingest\_watchlist.py},
    \texttt{scripts/ingest\_reviews.py}, the parser modules under
    \texttt{lib/parsers/}, and the central Supabase client in
    \texttt{lib/db.py}.

    \item \textbf{Analytical and machine-learning components:} To
    implement and evaluate a sequence of analytical modules for
    extracting competitive signals from the collected data. These
    modules include entrant/exit detection, listing and image-change
    detection, sponsored share-of-voice estimation, K-Means price-tier
    clustering, RoBERTa-based review sentiment classification, Brand
    Momentum Score (BMS), Listing Quality Score (LQS), rule-based alerts,
    and Prophet-based short-term price forecasting. In the codebase, this
    objective is implemented through \texttt{scripts/run\_analytics.py}
    and the individual \texttt{scripts/analyze\_*.py} modules, with
    model evaluation supported by \texttt{scripts/evaluate\_sentiment.py}
    and \texttt{scripts/evaluate\_forecast.py}.

    \item \textbf{Multi-channel intelligence delivery:} To deliver the
    computed intelligence through complementary interfaces: a Streamlit
    analytical dashboard for internal exploration, a lightweight
    Vercel-deployable web dashboard for compact market overviews, and an
    OpenClaw-powered Telegram layer for push-based briefs. The agentic
    layer consists of thirteen read-only Python skills under
    \texttt{openclaw/skills/}, corresponding OpenClaw wrappers under
    \texttt{openclaw/wrappers/}, and three Telegram-bound specialist
    agents: Competitor Spy, Sentiment Detective, and Momentum Strategist.
\end{enumerate}

\section{Research questions}

Guided by the objectives above, this thesis addresses the following
research questions:
\begin{enumerate}
    \item How can an actor-based data ingestion architecture be designed
    to collect and persist Amazon marketplace data reliably under the
    practical constraints of scheduled scraping, API cost, and schema
    consistency?

    \item To what extent can RoBERTa sentiment classification, Prophet
    price forecasting, K-Means price-tier clustering, perceptual hashing,
    and rule-based market metrics produce defensible competitive
    intelligence signals from short-window e-commerce data?

    \item How can a tool-using multi-agent delivery layer built on
    OpenClaw reduce the cognitive friction of dashboard-only monitoring
    by pushing concise, role-specific market briefs to Telegram?
\end{enumerate}

\section{Scope of the study}

The empirical scope of this thesis is limited to a 13-day pilot
deployment from April 17, 2026, to April 29, 2026. The market scope is
restricted to three Amazon consumer-electronics categories:
\textit{Gaming Keyboards}, \textit{True Wireless Earbuds}, and
\textit{Portable Chargers}. The system collects daily Top-50 category
rankings and detailed snapshots for a predefined watchlist of 30 ASINs,
while review ingestion is scheduled weekly because review extraction
consumes additional Apify credits. During the pilot window, the dataset
contains approximately 379 distinct ASINs observed across category
rankings and the fixed watchlist.

The technical scope includes the following implemented components:
\begin{itemize}
    \item an Apify-based acquisition layer using Amazon category and
    review actors;
    \item a Supabase PostgreSQL warehouse defined through
    version-controlled SQL migrations;
    \item Supabase Storage for product-image persistence and
    pHash-based image-change detection;
    \item Python ingestion, analytics, and evaluation scripts under
    \texttt{scripts/};
    \item GitHub Actions workflows for scheduled daily ingestion,
    weekly review collection, analytics-only reruns, and forecast-only
    reruns;
    \item a Streamlit dashboard under \texttt{dashboard/};
    \item a lightweight static web dashboard under
    \texttt{web-dashboard/}; and
    \item an OpenClaw delivery layer with thirteen read-only Python
    skills and three Telegram-bound specialist agents.
\end{itemize}

The study does not attempt to monitor the entire Amazon marketplace,
generalise findings to all product categories, or replace enterprise
commercial intelligence platforms. It also does not implement a custom
browser scraper, LangChain-based agent framework, Power BI dashboard, or
real-time streaming data system. All current Amazon data collection in
the implemented architecture flows through Apify, while all project
database access is centralised through \texttt{lib/db.py}.

\section{Structure of the thesis}

The remainder of this thesis is organised as follows:
\begin{itemize}
    \item \textbf{Chapter 2: Literature Review} reviews the literature
    on e-commerce data acquisition, competitive intelligence systems,
    machine-learning applications in retail analytics, dashboard-based
    decision support, and LLM-enabled multi-agent delivery systems.

    \item \textbf{Chapter 3: Methodology} presents the technical design
    of the proposed framework, including the data acquisition pipeline,
    database schema, analytical modules, machine-learning components,
    OpenClaw skill architecture, and evaluation methodology.

    \item \textbf{Chapter 4: Experimental Results} reports the empirical
    findings from the pilot deployment, including descriptive statistics,
    sentiment and forecasting evaluation, market-intelligence findings,
    alert-generation results, and operational observations from the
    multi-agent delivery layer.

    \item \textbf{Chapter 5: Conclusion and Future Work} summarises the
    main contributions, discusses the limitations of the current pilot
    implementation, and outlines future directions for longer-term data
    collection, richer modelling, and more advanced agentic delivery.
\end{itemize}
